%!TEX root =  main.tex

\chapter{Factor Graph Background}
In this chapter, we will first introduce factor graphs as a general tool for formulating estimation problems.  Following this, we will work through the specific example of classic discrete-time filtering/smoothing.  We will later show that continuous-time estimation is completely compatible with this common discrete-time setup, under specific conditions.  For a more thorough introduction to factor graphs, we refer the reader to~\cite{dellaertFactorGraphsRobot2017}.

\section{Factor Graphs}
\label{sec:fg}

Typically in state estimation, we are interested in computing a posterior density representing the likelihood of a state given some inputs and measurements of the form
\begin{equation}
p(\mbf{x} | \mbf{v}, \mbf{y}),
\end{equation}
where $\mbf{x} = (\mbf{x}_0, \ldots, \mbf{x}_K)$ is the state at times $t_0, \ldots, t_K$, $\mbf{v} = (\mbf{v}_0, \ldots, \mbf{v}_K)$ is a set of inputs/commands, and $\mbf{y} = (\mbf{y}_0, \ldots, \mbf{y}_K)$ is a set of measurements.  In most practical scenarios, we can factor this posterior density into a product of factors, each of which depends on only a subset of the variables in the problem.  For example, we might have
\begin{equation}
p(\mbf{x} | \mbf{v}, \mbf{y}) \propto \underbrace{p(\mbf{x}_0)}_{\phi(\mbf{x}_0)} \prod_{k=1}^K \underbrace{p(\mbf{x}_k | \mbf{x}_{k-1}, \mbf{v}_{k,k-1})}_{\phi(\mbf{x}_{k-1},\mbf{x}_k)} \prod_{k=0}^K \underbrace{p(\mbf{y}_k | \mbf{x}_k)}_{\varphi(\mbf{x}_k)}.
\end{equation}
where each $\phi(\cdot)$, is a {\em factor} of the posterior density.  The factors can be {\em unary} (involving one variable), {\em binary} (involving two variables), or higher-order (involving more than two variables).

The normalization constant associated with this factorization is typically expensive to compute so we often satisfy ourselves with finding the $\mbf{x}$ that maximizes the posterior density, or equivalently the $\mbf{x}$ that minimizes the negative log-likelihood cost:
\begin{equation}\label{eq:cost}
\mathcal{C}(\mbf{x}) = -\log \phi(\mbf{x}_0) - \sum_{k=1}^K \log \phi(\mbf{x}_k, \mbf{x}_{k-1}) - \sum_{k=0}^K \log \varphi(\mbf{x}_k).
\end{equation}
This is a common approach in estimation, referred to as \ac{MAP} estimation.

\begin{figure}[t]
     \includegraphics[width=\textwidth]{dtfg1.pdf}
     \caption{A simple factor graph for a classic estimation problem.  The unary factor on $\mbf{x}_0$ represents the prior knowledge about the initial state, the binary factors represent the motion model, and the unary factors represent the measurements.  In practice, we can have much more complicated factor graphs with loops.}
     \label{fig:dtfg1}
\end{figure}

We can visually understand the factors by assembling them into a {\em factor graph}, which is a bipartite graph with two types of nodes: variable nodes and factor nodes.  The variable nodes represent the variables in the problem, such as $\mbf{x}_0, \ldots, \mbf{x}_K$, and the factor nodes represent the factors, such as $\phi(\mbf{x}_0)$, $\phi(\mbf{x}_{k-1},\mbf{x}_k)$, and $\varphi(\mbf{x}_k)$.  The edges connect variable nodes to factor nodes, indicating which variables are involved in each factor.  This graphical representation allows us to see the dependencies between variables and factors in a very intuitive way.  Figure~\ref{fig:dtfg1} shows a simple factor graph for a classic discrete-time estimation problem.

\section{The Elimination Algorithm}

While it is possible to minimize the cost function in~\eqref{eq:cost} directly using matrix algebra, we here lean into the factor-graph representation and present the general {\em elimination algorithm}, an equivalent way to solve the minimization problem.  The key idea is to eliminate one variable at a time, updating the factors accordingly.  Once we have eliminated all variables, our factor graph has been transformed into another graph called a {\em Bayes net}; there is then a simple procedure to compute the desired posterior density from the Bayes net.

Mathematically, the elimination algorithm starts from a factor graph and ends with a new factorization of the posterior in the form
\begin{equation}\label{eq:elimfactors}
p(\mbf{x}) = p(\mbf{x}_0 | \mbf{s}_0) p(\mbf{x}_1 | \mbf{s}_1) \cdots p(\mbf{x}_K) = \prod_{k=0}^K p(\mbf{x}_k | \mbf{s}_k),
\end{equation}
where $\mbf{s}_k$ is the {\em separator} for $\mbf{x}_k$, defined as the set of variables on which $\mbf{x}_k$ is conditioned after it is eliminated.  Without loss of generality, here $k=0\ldots K$ is the order in which the variables have been eliminated. It is important to note that each variable's separator cannot include any variables that have already been eliminated, which means that once all the conditional densities are computed, we can work backwards from $K$ to $0$ to compute the posterior densities for each variable.  

Each step of the elimination algorithm consists of three operations:  (i) select a variable to eliminate, (ii) multiply all factors involving that variable together to form the {\em product factor}, and (iii) factor the product factor into a conditional density on the eliminated variable and a new factor on the separator variables.  The new factor is then added to the factor graph, and the eliminated variable is removed from the graph.  This process is repeated until all variables have been eliminated.  

In a bit more detail, consider the example factor graph on the left of Figure~\ref{fig:elim}.  Suppose we select $\mbf{x}_k$ as the variable to eliminate, then the product factor is given by
\begin{equation}
\psi(\mbf{x}_k, \mbf{s}_k) = \phi(\mbf{x}_k) \phi(\mbf{x}_k, \mbf{x}_i) \phi(\mbf{x}_k, \mbf{x}_j),
\end{equation}
where $\mbf{x}_i$ and $\mbf{x}_j$ are two other variables that share factors with $\mbf{x}_k$; the separator for $\mbf{x}_k$ is then $\mbf{s}_k = \{ \mbf{x}_i, \mbf{x}_j \}$.  The product factor is then factored into a conditional density on $\mbf{x}_k$ and a new factor on the separator variables:
\begin{equation}
\psi(\mbf{x}_k, \mbf{s}_k) = p(\mbf{x}_k | \mbf{s}_k) \, \phi(\mbf{s}_k).
\end{equation}
The key to this last factorization is that $\phi(\mbf{s}_k)$ has no dependence on $\mbf{x}_k$; it is in fact the factor that would result if $\mbf{x}_k$ were {\em marginalized} from factor graph. After all variables have been eliminated, we are left with a graph with only directed arrows called the Bayes net.

\begin{figure}[t]
     \includegraphics[width=\textwidth]{elim.pdf}
     \caption{Example of the elimination algorithm.  The left shows the factor graph before elimination, the middle shows the product factor formed by multiplying all factors involving $\mbf{x}_k$, and the right shows the result of eliminating $\mbf{x}_k$ from the factor graph, resulting in a new factor on the separator variables and a conditional density for $\mbf{x}_k$.}
     \label{fig:elim}
\end{figure}

The elimination algorithm is quite powerful in that it can be applied to any factor graph.  However, it is particularly useful for problems with a lot of structure, such as those arising in robotics, as we can exploit the sparsity of the factor graph in the various operations.  Here, we will only need to work with factors and densities that are Gaussian, and we use the next section to show that the elimination algorithm can be carried out in closed form when the factors are Gaussian.  This will allow us to carry out estimation efficiently and effectively.

\section{Gaussian Factors, Elimination, Fusion}
\label{sec:fgelimfus}

Now that we understand factor graphs and the general strategy for carrying out estimation, we show that when our factors are Gaussian we can carry out elimination in closed form.  We know that a Gaussian density has the form
\begin{equation}
p(\mbf{x}) \propto \exp\left( -\frac{1}{2} (\mbf{x} - \mbs{\mu})^T \mbs{\Sigma}^{-1} (\mbf{x} - \mbs{\mu}) \right) = \exp\left( -\frac{1}{2} || \mbf{x} - \mbs{\mu} ||^2_\mbs{\Sigma} \right).
\end{equation}
After taking the negative logarithm, we will only need the $|| \mbf{x} - \mbs{\mu} ||^2_\mbs{\Sigma}$ part, which forms a term in the negative-log likelihood cost that we are seeking to minimize; we will often then label our factors on our factor graphs using this cost term for simplicity.  We may also sometimes write $p(\mbf{x}) \sim \mathcal{N}(\mbs{\mu}, \mbs{\Sigma})$ as a shorthand for Gaussian densities.

\begin{figure}[t]
     \includegraphics[width=\textwidth]{gausselim.pdf}
     \caption{With Gaussian factors, eliminating $\mbf{x}$ results in a new Gaussian factor and a Gaussian conditional density for $\mbf{x}$.}
     \label{fig:gausselim}
\end{figure}

Consider the situation on the left of Figure~\ref{fig:gausselim}; we have $\mbf{x}$ linked to two other variables via Gaussian factors (i.e., $\mbf{x}$'s separator).  When we eliminate $\mbf{x}$, this results in the situation on the right, with a new binary factor between the separator variables as well as a Gaussian conditional density for $\mbf{x}$.  Note, we assume that $\mbf{u}_1$ depends on $\mbf{x}_1$ and $\mbf{u}_2$ depends on $\mbf{x}_2$, typically linearly, but the elimination procedure still works if $\mbf{x}$ has a binary factor on one side and a unary factor on the other. 

The Gaussian elimination in Figure~\ref{fig:gausselim} follows from the very handy sum-of-squares identity,
\begin{equation}\label{eq:sumofsquares}
|| \mbf{u}_1 - \mbf{A}_1 \mbf{x} ||^2_{\mbf{Q}_1} + || \mbf{u}_2 - \mbf{A}_2 \mbf{x} ||^2_{\mbf{Q}_2}  = || \mbf{A}_1^{-1} \mbf{u}_1 - \mbf{A}_2^{-1} \mbf{u}_2 ||^2_{\mbf{Q}_{12}} + || \mbf{x} - \mbs{\mu}||^2_\mbs{\Sigma},
\end{equation}
where $\mbf{Q}_{12}$, $\mbs{\Sigma}$, $\mbs{\mu}$,  are defined in Figure~\ref{fig:gausselim} and $\mbf{A}_1$ and $\mbf{A}_2$ are for now assumed invertible.  We will use this closed-form Gaussian elimination time and again in what follows.  

\begin{figure}[t]
     \centering
     \includegraphics[width=0.8\textwidth]{gaussfusion.pdf}
     \caption{We can fuse multiple Gaussian factors into a single one in order to declutter our factor graph.}
     \label{fig:gaussfusion}
\end{figure}

Another operation that we will perform frequently with Gaussian factors is fusion.  By this we mean lumping together multiple factors into a single one as depicted in Figure~\ref{fig:gaussfusion} for unary factors.  In fact, it is a well-known result that we can fuse an arbitrary number of Gaussians in this way, allowing us to declutter our factor graph.

We are now prepared to discuss classic discrete-time estimation from a factor-graph perspective.

\section{Linear Discrete-Time Estimation}
\label{sec:ldte}

Classic discrete-time state estimation has a chain-like factor graph as depicted in Figure~\ref{fig:dtfg1}.  There is a unary initial-condition factor capturing what we know about the state at the initial time, binary motion-model factors predicting how the system moves from one discrete time to the next, and unary measurement factors that tell us something about the state at each discrete time. 

% \begin{figure}[b]
%      \centering
%      \includegraphics[width=0.95\textwidth]{dtfg.pdf}
%      \caption{Factor graph for classic discrete-time filtering/smoothing estimators.}
%      \label{fig:dtfg}
% \end{figure}

For example, we might have a stochastic linear system of the form
\beqn{linsys}
\mbf{x}_k & = & \mbf{A}_{k,k-1} \mbf{x}_{k-1} + \mbf{v}_{k,k-1} + \mbf{w}_{k,k-1},  \\
\mbf{y}_k & = & \mbf{C}_k \mbf{x}_k + \mbf{n}_k,
\eeqn
where $k \in 0 \ldots K$ is the discrete-time index, $\mbf{x}_k$ is the state, $\mbf{v}_{k,k-1}$ is an input/command, $\mbf{y}_k$ is an output/measurement, $\mbf{w}_{k,k-1} \sim \mathcal{N}(\mbf{0}, \mbf{Q}_{k,k-1})$ is Gaussian process noise, and $\mbf{n}_k \sim \mathcal{N}(\mbf{0}, \mbf{R}_k)$ is Gaussian measurement noise.  We additionally assume that at $k=0$ the initial condition is drawn from a Gaussian, $\mbf{x}_0 \sim \mathcal{N}\left(\pri{\mbf{x}}_0, \pri{\mbf{P}}_0\right)$.  In this case, we can populate our factor graph as in Figure~\ref{fig:dtfg2}, where all of the factors are now Gaussian.  

An important point to make when looking at Figure~\ref{fig:dtfg2} is that the factor graph has a lot of structure in the form of sparsity.  Not all nodes are connected to all other nodes; rather, each node is connected only to its temporal neighbours.  This sparsity arises because $\mbf{x}$ is a {\em Markovian} state for~\eqref{eq:linsys}.  In other words, it captures everything we need to know to predict the next state of the system.  

Our state-estimation problem is now to use the initial-state knowledge, $\pri{\mbf{x}}_0$, the inputs, $\mbf{v}_{k,k-1}$, and the measurements, $\mbf{y}_k$, to work out the likelihood of $\mbf{x}$ taking on all possible values.   This can be done exactly for this linear-Gaussian system; later for nonlinear systems we will require some approximations.  We will show that by using a combination of the Gaussian elimination and fusion operations from Section~\ref{sec:fgelimfus}, we can recover the classic \ac{RTS} smoother, the forward pass of which is the famous \ac{KF}.

\begin{figure}[t]
     \centering
     \includegraphics[width=\textwidth]{dtfg2.pdf}
     \caption{Populated factor graph for the discrete-time estimation of the linear system in~\eqref{eq:linsys}.}
     \label{fig:dtfg2}
\end{figure}

Looking to Figure~\ref{fig:dtfg2}, our estimation game plan will be to first fuse the two unary factors on $\mbf{x}_0$, then eliminate $\mbf{x}_0$, then repeat this process with $\mbf{x}_1$ and so on until we arrive at $\mbf{x}_K$.  The result of this forward pass will be a Bayes net.  After this we will carry out a backward pass to work out the posterior densities for $\mbf{x}_K$ all the way back to $\mbf{x}_0$.

\begin{figure}[p]
     \centering
     \includegraphics[width=\textwidth]{dtfg3.pdf}
     \caption{Result after fusing the two unary factors on $\mbf{x}_0$ in Figure~\ref{fig:dtfg2}.}
     \label{fig:dtfg3}
\end{figure}
\begin{figure}[p]
     \centering
     \includegraphics[width=\textwidth]{dtfg4.pdf}
     \caption{Result after eliminating $\mbf{x}_0$ from the factor graph in Figure~\ref{fig:dtfg3}.}
     \label{fig:dtfg4}
\end{figure}
\begin{figure}[p]
     \centering
     \hspace*{0.5in}\includegraphics[width=0.87\textwidth]{dtfg5.pdf}
     \caption{Result after fusing the two unary factors on $\mbf{x}_1$ in Figure~\ref{fig:dtfg4} and then also eliminating $\mbf{x}_1$.}
     \label{fig:dtfg5}
\end{figure}

Figure~\ref{fig:dtfg3} shows the result of fusing the two unary factors on $\mbf{x}_0$ into a single Gaussian factor.  Figure~\ref{fig:dtfg4} then shows the result of eliminating $\mbf{x}_0$ from the factor graph, resulting in a new unary factor on $\mbf{x}_1$ as well as the conditional density for $\mbf{x}_0$.  We now notice that the remaining factor graph looks identical to the one we started with, only shifted forward by one time step.  Figure~\ref{fig:dtfg5} shows what happens after another round of fusing the unary factors on $\mbf{x}_1$ then eliminating $\mbf{x}_1$.  If we keep applying this same procedure all the way to time step $K$, our factor graph is completely eliminated and we are left with the Bayes net in Figure~\ref{fig:dtfg6}.  We can then work backward from $\mbf{x}_K$ back to $\mbf{x}_0$ using the conditional densities to compute the posterior densities for each state.

\begin{figure}[t]
     \centering
     \includegraphics[width=0.85\textwidth]{dtfg6.pdf}
     \caption{After eliminating all the variables in a forward pass, we are left with a Bayes net.  We can then carry out a backward pass to calculate the posterior densities on all the variables.}
     \label{fig:dtfg6}
\end{figure}

After some algebraic manipulations, it can be shown that this procedure is identical to the classic \ac{RTS} smoother, consisting of the same forward and backward passes:
\begin{subequations}
\label{eq:rts}
\begin{eqnarray}
\mbox{forward:} \quad\;\;\; & & \nonumber \\
\mbox{($k=1\ldots K$)} \; & & \nonumber \\
\pri{\mbf{P}}_{k,f} & = & \mbf{A}_{k,k-1} \est{\mbf{P}}_{k-1,f} \mbf{A}_{k,k-1}^T + \mbf{Q}_{k,k-1}, \\
\pri{\mbf{x}}_{k,f} & = & \mbf{A}_{k,k-1} \est{\mbf{x}}_{k-1,f} + \mbf{v}_{k,k-1}, \\
\mbf{K}_k & = & \pri{\mbf{P}}_{k,f} \mbf{C}_k^T \left( \mbf{C}_k \pri{\mbf{P}}_{k,f} \mbf{C}_k^T + \mbf{R}_k \right)^{-1}, \\
\est{\mbf{P}}_{k,f} & = & \left( \mbf{1} - \mbf{K}_k \mbf{C}_k \right) \pri{\mbf{P}}_{k,f}, \\
\est{\mbf{x}}_{k,f} & = & \pri{\mbf{x}}_{k,f} + \mbf{K}_k \left( \mbf{y}_k - \mbf{C}_k \pri{\mbf{x}}_{k,f} \right), \\
\mbox{backward:}\quad & & \nonumber \\
\mbox{($k=K\ldots 1$)} \; & & \nonumber \\
\est{\mbf{x}}_{k-1} & = & \est{\mbf{x}}_{k-1,f} + \left(\est{\mbf{P}}_{k-1,f} \mbf{A}_{k,k-1}^T \pri{\mbf{P}}_{k,f}^{-1} \right)\left( \est{\mbf{x}}_k - \pri{\mbf{x}}_{k,f} \right), \qquad\quad \\
\est{\mbf{P}}_{k-1} & = & \est{\mbf{P}}_{k-1,f} + \left( \est{\mbf{P}}_{k-1,f} \mbf{A}_{k,k-1}^T \pri{\mbf{P}}_{k,f}^{-1}\right) \left( \est{\mbf{P}}_k - \pri{\mbf{P}}_{k,f} \right) \nonumber \\  &  & \qquad\qquad\qquad\quad \times \; \left( \est{\mbf{P}}_{k-1,f} \mbf{A}_{k,k-1}^T \pri{\mbf{P}}_{k,f}^{-1}\right)^T \label{eq:rtsbackcov3}, 
\end{eqnarray}
\end{subequations}
which are initialized with 
\begin{subequations}
\begin{eqnarray}
\est{\mbf{P}}_{0,f} & = & ( \mbf{1} - \mbf{K}_0 \mbf{C}_0 ) \pri{\mbf{P}}_0, \\
\est{\mbf{x}}_{0,f} & = & \pri{\mbf{x}}_0 + \mbf{K}_0 ( \mbf{y}_0 - \mbf{C}_0 \pri{\mbf{x}}_0), \\
\est{\mbf{x}}_K & = & \est{\mbf{x}}_{K,f}, \\
\est{\mbf{P}}_K & = & \est{\mbf{P}}_{K,f},
\end{eqnarray}
\end{subequations}
and $\mbf{K}_0 = \pri{\mbf{P}}_{0} \mbf{C}_0^T \left( \mbf{C}_0 \pri{\mbf{P}}_{0} \mbf{C}_0^T + \mbf{R}_0 \right)^{-1}$.  We will also be interested in computing the cross-covariance between two consecutive states, which is given by
\begin{equation}
\est{\mbf{P}}_{k,k-1} = \est{\mbf{P}}_{k} \left( \est{\mbf{P}}_{k-1,f} \mbf{A}_{k,k-1}^T \pri{\mbf{P}}_{k,f}^{-1}\right)^T
\end{equation}
and can be calculated during the backward pass.  The computational complexity of the \ac{RTS} smoother is $O(K)$, in other words, linear in the length of the trajecotry.

The forward pass is the very famous \acf{KF}, which we now understand from a factor-graph perspective.  The $p(\mbf{x}_k | \mbf{v}_{0:k}, \mbf{y}_{0:k}) = \mathcal{N}\left(\est{\mbf{x}}_{k,f}, \est{\mbf{P}}_{k,f} \right)$ quantities can be referred to as {\em filtered} (accounting only for up-to-the-present measurements) while the $p(\mbf{x}_k | \mbf{v}_{0:K}, \mbf{y}_{0:K}) = \mathcal{N}\left( \est{\mbf{x}}_k, \est{\mbf{P}}_k \right)$ quantities are {\em smoothed} (accounting for all measurements: past, present, and future); see Figure~\ref{fig:filtersmoother}.

\begin{figure}[b]
     \centering
     \includegraphics[width=0.95\textwidth]{filtersmoother.pdf}
     \caption{Contrasting filtering vs. smoothing.  We are primarily interested in smoothing in this article.}
     \label{fig:filtersmoother}
\end{figure}

It turns out that this discrete-time estimation setup is completely compatible with extending our results to continuous time, we just need to fill in the details of~\eqref{eq:linsys} in the right way.  This will be the subject of the next chapter.


